% 定性示意，非实验数据；横坐标表示压缩后的训练进程。
\begin{tikzpicture}[x=1mm,y=1mm,font=\kaishu\footnotesize]
 \fill[BookPaper] (19,0) rectangle (65,40);
 \draw[book arrow] (0,0)--(99,0) node[below left]{训练进程};
 \draw[book arrow] (0,0)--(0,46) node[above]{错误比例};
 \draw[BookMuted,dashed] (19,0)--(19,40);
 \draw[BookMuted,dashed] (65,0)--(65,40);
 \draw[BookTeal,very thick,dashed] (2,36) .. controls (10,31) and (12,4) .. (19,2)
  .. controls (36,1) and (70,1) .. (96,1);
 \draw[BookAmber,very thick] (2,37) .. controls (23,36) and (50,36) .. (60,34)
  .. controls (68,30) and (69,6) .. (79,3)
  .. controls (85,2) and (91,2) .. (96,2);
 \node[text=BookAmber,font=\scriptsize] at (42,40){测试错误率};
 \node[text=BookTeal,font=\scriptsize] at (41,5){经验零一风险};
 \node[text=BookMuted,font=\scriptsize,align=center] at (42,20){样本已被正确分类\\新输入上的错误仍然较多};
 \node[below,font=\scriptsize,align=center] at (19,0){完成分类拟合};
 \node[below,font=\scriptsize,align=center] at (70,-6){延迟出现的泛化改善};
\end{tikzpicture}
